% TOP_PROBLEM: {"id": 30006970, "title": "Ruzsa-Szemeredi (6,3) Extremal-Order Problem", "category_id": 2, "category": {"id": 2, "name": "combinatorics", "display_name": "Combinatorics", "description": "Counting problems, graph theory, discrete structures.", "slug": "combinatorics", "order_index": 2, "created_at": "2026-07-31T15:26:25.671Z"}, "status": "open"}
% FIRST_POSTED: 2026-09-25
% ENTRY_KIND: statement_only
% !TeX program = lualatex
% Definition and sources only; this file is not an LLM solution attempt.
\documentclass[11pt]{article}
\usepackage[margin=25mm]{geometry}
\usepackage{fontspec}
\setmainfont{Latin Modern Roman}
\usepackage{amsmath,amssymb,mathtools}
\usepackage{unicode-math}
\setmathfont{Latin Modern Math}
\usepackage{xurl,hyperref}
\hypersetup{colorlinks=true,urlcolor=blue}
\setcounter{secnumdepth}{0}
\setlength{\parindent}{0pt}
\setlength{\parskip}{5pt}
\setlength{\emergencystretch}{3em}
\begin{document}
\section*{\texorpdfstring{Ruzsa-Szemeredi (6,3) Extremal-Order Problem}{Ruzsa-Szemeredi (6,3) Extremal-Order Problem}}\label{30006970}
\subsection{Definitions and mathematical statement}
A simple three-uniform hypergraph on an $n$-element vertex set has distinct three-element subsets as edges. Define
\[
 f(n,6,3)=\max\{|E(H)|:\ |V(H)|=n,
 \ \forall U\subseteq V(H),\ |U|\le6,
 \ \#\{e\in E(H):e\subseteq U\}\le2\}.
\]
Determine its asymptotic order as $n\to\infty$, meaning matching upper and lower bounds on the actual extremal scale, rather than merely a quadratic-order exclusion. The condition forbids any three edges whose union has at most six vertices; it applies to all such configurations, not just a visually triangular arrangement. No multiplicities are allowed among edges.
\par\smallskip\textit{Formulation and concept references:} \href{https://arxiv.org/html/2003.02754}{[S1]}.

\subsection{Short English statement}
How many triples can a large vertex set support when no collection of at most six vertices contains three of them?
\subsection{Sources}
\begin{itemize}
\item[S1]Gowers and Janzer, Generalizations of the Ruzsa-Szemeredi and rainbow Turan problems for cliques (2020).
\url{https://arxiv.org/html/2003.02754}.
\textit{Formulation source.}
\item[S2]Status source for Ruzsa-Szemeredi (6,3) Extremal-Order Problem.
\url{https://people.math.ethz.ch/~sudakovb/Brown-Erdos-Sos-power-saving.pdf}.
\textit{Further reference; not a proof endorsement.}
\end{itemize}
\end{document}
